\chapter{Практическая реализация фреймворка для микроконтроллера MIK32}

\label{cha:framework}

\section{Архитектура фреймворка}

Архитектура разрабатываемого фреймворка представляет собой многоуровневую структуру, обеспечивающую четкое разделение между аппаратно-зависимыми и аппаратно-независимыми компонентами. Данная архитектура позволяет создать гибкую и масштабируемую платформу для разработки серверных компонентов, обеспечивая возможность переноса решений между различными аппаратными конфигурациями~\cite{mikron_mik32_amur}.

Фреймворк построен на основе абстракции, которая обеспечивает четкое разделение функциональности между различными компонентами системы. Такой подход соответствует современным практикам разработки встраиваемых систем, где необходимо обеспечить гибкость, масштабируемость и переносимость программного обеспечения~\cite{embedded_abstraction_layers}.

В контексте разрабатываемого фреймворка архитектура разделяется на следующие основные уровни:

\begin{itemize}
	\item \textbf{Архитектурно-зависимый уровень (arch)} -- отвечает за взаимодействие с аппаратными ресурсами микроконтроллера MIK32, включая прямое управление регистрами и периферийными устройствами
	\item \textbf{Платформозависимый уровень (platform)} -- обеспечивает абстракцию конкретной аппаратной платформы, адаптируя архитектурно-зависимые компоненты к особенностям конкретной платы
	\item \textbf{Общий уровень (включает src и include)} -- содержит высокоуровневые компоненты, которые зависят от платформы, но не являются частью платформозависимой архитектуры
	\item \textbf{Вспомогательные инструменты (tools)} -- обеспечивает инструментарий для работы с микроконтроллером, включая системы сборки и отладочные утилиты
\end{itemize}

Такая стратификация архитектуры позволяет изолировать изменения, связанные с аппаратной платформой, и минимизировать их влияние на остальные компоненты системы. Это соответствует принципу разделения обязанностей (separation of concerns), что делает фреймворк более гибким и адаптируемым к различным условиям применения~\cite{layered_architecture_embedded}. На рисунке~\ref{fig:struct_diagr} показано разделение данных уровней.

\begin{figure}[H]
	\centering
	\includegraphics[width=0.5\textwidth]{inc/img/struct-diagr}
	\caption{Архитектура фреймворка}
	\label{fig:struct_diagr}
\end{figure}

Представленная архитектурная диаграмма демонстрирует иерархическое построение фреймворка, где каждый уровень обеспечивает определенную степень абстракции. Нижний уровень взаимодействует непосредственно с аппаратным обеспечением микроконтроллера, средние уровни предоставляют платформо-специфичные адаптации, а верхние уровни содержат переносимую бизнес-логику приложений.

\subsection{Архитектурно-зависимая часть (arch)}

Архитектурно-зависимая часть фреймворка представляет собой фундаментальный уровень взаимодействия с аппаратным обеспечением микроконтроллера MIK32. Данный уровень обеспечивает абстракцию аппаратных особенностей микроконтроллера, предоставляя вышестоящим уровням унифицированный интерфейс доступа к функциональности.

Архитектурно-зависимая часть содержит следующие компоненты:

\begin{itemize}
	\item \textbf{src (исходный код)} -- реализация низкоуровневых драйверов
	\item \textbf{include (заголовочные файлы)} -- определения интерфейсов и структур данных
	\item \textbf{lib (библиотеки)} -- повторно используемые компоненты
	\item \textbf{framework-sdk} -- библиотеки HAL и LL уровней абстракции
\end{itemize}

В каталоге \texttt{src} размещается исходный код, специфичный для архитектуры микроконтроллера MIK32. Здесь содержатся реализации низкоуровневых драйверов, обеспечивающих прямое взаимодействие с аппаратными ресурсами. Эти драйверы инкапсулируют детали работы с регистрами и прерываниями микроконтроллера, предоставляя более высоким уровням абстракции чистый программный интерфейс.

Архитектурно-зависимый исходный код отвечает за следующие функции:

\begin{itemize}
	\item Инициализация и настройка периферийных устройств согласно техническим характеристикам MIK32~\cite{mik32_technical_spec}
	\item Управление энергопотреблением с использованием доменов питания микроконтроллера
	\item Обработка аппаратных прерываний через программируемый контроллер прерываний
	\item Взаимодействие с регистрами микроконтроллера через карту памяти
	\item Реализация базовых операций ввода-вывода
\end{itemize}

Важно отметить, что код в этом каталоге напрямую зависит от архитектуры микроконтроллера MIK32 и не предназначен для непосредственного использования вышестоящими уровнями абстракции.

Каталог \texttt{framework-sdk} представляет собой наиболее важный компонент архитектурно-зависимой части фреймворка. Он содержит библиотеки HAL (Hardware Abstraction Layer) и LL (Low Level), которые предоставляют различные уровни абстракции для взаимодействия с аппаратным обеспечением микроконтроллера MIK32~\cite{stm32_hal_ll_comparison}.

LL (Low Level) библиотека предоставляет низкоуровневый доступ к аппаратным ресурсам микроконтроллера. Она спроектирована для обеспечения максимальной производительности и минимальных накладных расходов при работе с периферийными устройствами. HAL (Hardware Abstraction Layer) библиотека предоставляет более высокий уровень абстракции, скрывая детали реализации аппаратной платформы за унифицированным интерфейсом~\cite{stm32_hal_ll_comparison}.

\subsection{Платформозависимая часть (platform)}

Платформозависимая часть фреймворка обеспечивает абстракцию конкретной аппаратной платформы, на которой работает микроконтроллер MIK32. В отличие от архитектурно-зависимой части, которая фокусируется на особенностях самого микроконтроллера, платформозависимая часть адресует специфику конкретной платы или устройства, на котором установлен микроконтроллер.

Платформозависимая часть содержит следующие компоненты:

\begin{itemize}
	\item \textbf{src (исходный код)} -- адаптация к особенностям конкретной платы
	\item \textbf{include (заголовочные файлы)} -- платформо-специфичные определения
\end{itemize}

Платформозависимый исходный код отвечает за следующие функции:

\begin{itemize}
	\item Инициализация платформо-специфичных периферийных устройств
	\item Конфигурация тактовых сигналов с учетом особенностей платформы
	\item Управление питанием с учетом схемотехники конкретной платы
	\item Настройка и управление внешними компонентами (дисплеи, сенсоры и т.д.)
	\item Реализация механизмов взаимодействия с внешними устройствами
\end{itemize}

\subsection{Система сборки}

Файл \texttt{Makefile} управляет процессом сборки фреймворка и приложений на его основе. Он обеспечивает интеграцию всех компонентов архитектуры и автоматизирует процесс компиляции, линковки и генерации исполняемого кода. Особенностью системы сборки в данном фреймворке является использование символьных ссылок, что обеспечивает гибкость и масштабируемость процесса сборки.

\section{Модуль GPIO}

\subsection{Архитектура и основные компоненты}

Подсистема GPIO микроконтроллера MIK32 реализована на основе модульной архитектуры, что обеспечивает гибкость и расширяемость при работе с различными аппаратными контроллерами портов ввода-вывода. Архитектурное решение основано на принципах, применяемых в подсистеме GPIO ядра Linux, адаптированных для встраиваемых систем~\cite{linux_gpio_subsystem}.

Центральным элементом системы выступает структура \texttt{gpio\_chip}, которая инкапсулирует все необходимые сведения о конкретном GPIO-контроллере. Она содержит имя чипа, указатель на приватные данные, таблицу операций, базовый номер выводов, количество доступных линий и массив дескрипторов этих линий. Такой подход позволяет абстрагироваться от специфики низкоуровневой реализации и использовать унифицированный API для работы с различными типами GPIO-контроллеров, включая как встроенные порты микроконтроллера, так и внешние расширители, такие как PCA9554.

Структура дескриптора чипа GPIO представлена в на рисунке\ref{fig:struct_gpio}. Данная структура является основой для реализации полиморфизма в подсистеме GPIO, позволяя использовать единый интерфейс для работы с различными типами контроллеров.

\begin{figure}[H]
	\centering
	\includegraphics[width=0.9\textwidth]{inc/img/struct_gpio}
	\caption{Структура дескриптора чипа GPIO}
	\label{fig:struct_gpio}
\end{figure}


\begin{listing}[H]
	\cfile{inc/src/gpio-chip-struct.c}
	\caption{Програмная реализация структуры дескриптора чипа GPIO}
	\label{lst:gpio_chip}
\end{listing}

Листинг~\ref{lst:gpio_chip} демонстрирует архитектурный паттерн, где структура \texttt{gpio\_chip} содержит указатель на таблицу виртуальных функций \texttt{gpio\_chip\_ops}. Это решение обеспечивает возможность динамического связывания операций в зависимости от типа GPIO-контроллера. Поле \texttt{priv\_data} позволяет каждому драйверу сохранять свои специфические данные, а массив \texttt{descs} предоставляет быстрый доступ к дескрипторам отдельных выводов.

Каждый отдельный вывод описывается структурой \texttt{gpio\_desc}, представленной в листинге~\ref{lst:gpio_desc}. Эта структура обеспечивает детализированное управление состоянием каждого GPIO вывода.

\begin{listing}[H]
	\cfile{inc/src/gpio-desc-struct.c}
	\caption{Структура дескриптора GPIO вывода}
	\label{lst:gpio_desc}
\end{listing}

Структура \texttt{gpio\_desc} в листинге~\ref{lst:gpio_desc} инкапсулирует все необходимые атрибуты отдельного GPIO вывода. Битовые флаги \texttt{flags} позволяют гибко настраивать поведение вывода.

\subsection{Система флагов конфигурации}

Гибкость конфигурирования обеспечивается за счёт использования набора битовых флагов, каждый из которых отвечает за отдельный аспект поведения вывода. Система флагов, представленная в листинге~\ref{lst:gpio_flags}, позволяет задать направление (вход/выход), начальное состояние (логическая единица или ноль), тип выхода (push-pull или open-drain), наличие подтягивающих резисторов (pull-up/pull-down), полярность (active-high/active-low), а также назначение вывода для работы с прерываниями и другие параметры.

\begin{listing}[H]
	\cfile{inc/src/gpio-flags.c}
	\caption{Флаги конфигурации GPIO выводов}
	\label{lst:gpio_flags}
\end{listing}

Использование битовых флагов в листинге~\ref{lst:gpio_flags} представляет собой решение для хранения конфигурационной информации GPIO выводов. Такой подход минимизирует потребление памяти и обеспечивает быстрые операции над множественными атрибутами вывода. Флаги могут комбинироваться с помощью битовых операций, что позволяет создавать сложные конфигурации с минимальными вычислительными затратами.

\subsection{Управление жизненным циклом GPIO ресурсов}

GPIO-подсистема предусматривает динамическое добавление и удаление контроллеров, что реализовано через связанный список всех зарегистрированных чипов. Процесс регистрации нового GPIO-контроллера реализован в функции \texttt{gpiochip\_add}, код которой представлен в листинге~\ref{lst:gpiochip_add}.

\begin{listing}[H]
	\cfile{inc/src/gpiochip-add.c}
	\caption{Функция регистрации GPIO чипа}
	\label{lst:gpiochip_add}
\end{listing}

Функция \texttt{gpiochip\_add} в листинге~\ref{lst:gpiochip_add} выполняет комплексную инициализацию нового GPIO-контроллера в системе. Алгоритм автоматически подбирает свободный диапазон номеров для GPIO линий нового контроллера, инициализирует дескрипторы всех доступных выводов и интегрирует контроллер в глобальную структуру данных подсистемы. Проверка корректности параметров и обработка ошибок обеспечивают надежность работы системы.

Для управления доступом к отдельным выводам реализован механизм захвата и освобождения ресурсов. Функции управления ресурсами GPIO представлены в листинге~\ref{lst:gpio_get_put}.

\begin{listing}[H]
	\cfile{inc/src/gpio-get-put.c}
	\caption{Функции захвата и освобождения GPIO выводов}
	\label{lst:gpio_get_put}
\end{listing}

Механизм управления ресурсами в листинге~\ref{lst:gpio_get_put} предотвращает конфликты доступа к GPIO выводам между различными компонентами системы. Функция \texttt{gpiod\_get} проверяет доступность запрашиваемого вывода и резервирует его для конкретного потребителя, в то время как \texttt{gpiod\_put} освобождает ресурс для повторного использования. Такой подход обеспечивает детерминированное поведение системы и упрощает отладку конфликтов ресурсов.

\section{Модуль I2C}

Микроконтроллер MIK32 включает в себя два контроллера I2C, которые обеспечивают взаимодействие с различными периферийными устройствами~\cite{mik32_technical_spec}. Подсистема I2C является критически важным компонентом для обеспечения связи между микроконтроллером и различными периферийными устройствами, такими как сенсоры, дисплеи и другие интегральные схемы.

\subsection{Архитектура подсистемы I2C}

Разработанная подсистема I2C для фреймворка MIK32 представляет собой многоуровневую архитектуру, обеспечивающую абстракцию от аппаратных особенностей контроллера. Архитектурный подход основан на принципах, применяемых в подсистеме I2C ядра Linux~\cite{linux_i2c_subsystem}, адаптированных для требований встраиваемых систем.

Основу архитектуры составляют три ключевые структуры данных: адаптер шины I2C, клиентское устройство и сообщение для передачи данных. Такой подход позволяет создать универсальный интерфейс для работы с различными устройствами на шине I2C, независимо от конкретной реализации аппаратного контроллера.

Архитектура построена по принципу разделения ответственности, где каждый компонент выполняет строго определенные функции:

\begin{itemize}
	\item \textbf{Адаптер шины} (\texttt{struct i2c\_adapter}) -- представляет собой абстракцию физического контроллера I2C и управляет всеми операциями на конкретной шине
	\item \textbf{Клиентские устройства} (\texttt{struct i2c\_client}) -- представляют периферийные устройства, подключенные к шине, и содержат информацию об их адресах и режимах работы
	\item \textbf{Сообщения} (\texttt{struct i2c\_msg}) -- инкапсулируют данные для передачи и определяют параметры транзакций
\end{itemize}

Фреймворк поддерживает стандартные режимы работы I2C, включая Standard Mode (100 кГц), Fast Mode (400 кГц), Fast Mode Plus (1 МГц) и High Speed Mode (3.4 МГц)~\cite{i2c_specification}, что обеспечивает совместимость с широким спектром периферийных устройств.

\subsection{Структуры данных и константы}

Подсистема I2C использует набор предопределенных констант для различных скоростей работы шины, как показано в листинге~\ref{lst:i2c_speeds}.

\begin{listing}[H]
	\cfile{inc/src/i2c-speeds.c}
	\caption{Константы скоростей I2C}
	\label{lst:i2c_speeds}
\end{listing}

Константы скоростей в листинге~\ref{lst:i2c_speeds} определяют стандартные режимы работы шины I2C согласно официальной спецификации. Использование именованных констант повышает читаемость кода и предотвращает ошибки, связанные с использованием магических чисел. Поддержка различных скоростей позволяет оптимизировать производительность системы в зависимости от требований конкретных периферийных устройств.

Центральной структурой подсистемы является дескриптор адаптера шины I2C, представленный на рисунке~\ref{fig:struct_i2c}.


\begin{figure}[H]
	\centering
	\includegraphics[width=0.9\textwidth]{inc/img/struct_i2c}
	\caption{Структура дескриптора чипа GPIO}
	\label{fig:struct_i2c}
\end{figure}


\begin{listing}[H]
	\cfile{inc/src/i2c-adapter.c}
	\caption{Програмная реализация дескриптора адаптера I2C}
	\label{lst:i2c_adapter}
\end{listing}

Структура \texttt{i2c\_adapter} в листинге~\ref{lst:i2c_adapter} инкапсулирует все параметры, необходимые для управления шиной I2C. Список клиентских устройств \texttt{client\_list} упрощает управление подключенными устройствами.
Процесс регистрации нового адаптера I2C в системе выполняется функцией, представленной в листинге~\ref{lst:i2c_adapter_add}.

\begin{listing}[H]
	\cfile{inc/src/i2c-adapter-add.c}
	\caption{Функция добавления адаптера I2C}
	\label{lst:i2c_adapter_add}
\end{listing}

Функция \texttt{i2c\_adapter\_add} в листинге~\ref{lst:i2c_adapter_add} выполняет комплексную инициализацию адаптера I2C в системе. Алгоритм автоматически присваивает уникальный номер адаптеру, инициализирует связанный список клиентских устройств и интегрирует адаптер в глобальную структуру данных подсистемы I2C. Проверка корректности параметров обеспечивает надежность работы системы.

\subsection{Работа с клиентскими устройствами}

Клиентские устройства на шине I2C представляются через специальную структуру, содержащую информацию об адресе устройства, ссылку на родительский адаптер и callback-функцию для обработки событий в режиме ведомого. Процесс регистрации клиентского устройства представлен в листинге~\ref{lst:i2c_client_register}.

\begin{listing}[H]
	\cfile{inc/src/i2c-client-register.c}
	\caption{Функция регистрации клиентского устройства}
	\label{lst:i2c_client_register}
\end{listing}

Функция \texttt{i2c\_client\_register} в листинге~\ref{lst:i2c_client_register} обеспечивает интеграцию нового клиентского устройства в подсистему I2C. Проверка уникальности адреса предотвращает конфликты на шине, а добавление клиента в список адаптера упрощает управление подключенными устройствами. Callback-функция позволяет устройству реагировать на события шины в режиме ведомого.

\subsection{Передача данных в режиме ведущего}

Передача данных в режиме ведущего реализована через функцию \texttt{i2c\_master\_xfer}, представленную в листинге~\ref{lst:i2c_master_xfer}, которая принимает массив сообщений для передачи и обеспечивает их последовательную обработку.

\begin{listing}[H]
	\cfile{inc/src/i2c-master-xfer.c}
	\caption{Функция передачи данных в режиме ведущего}
	\label{lst:i2c_master_xfer}
\end{listing}

Функция \texttt{i2c\_master\_xfer} в листинге~\ref{lst:i2c_master_xfer} реализует универсальный механизм передачи данных по шине I2C. Использование массива сообщений позволяет выполнять сложные транзакции, включающие множественные операции чтения и записи в рамках одной атомарной операции.

\subsection{Обработка прерываний и событий}

Подсистема I2C включает механизмы обработки прерываний и событий, которые позволяют управлять передачей данных без блокировки процессора. Система прерываний построена по иерархическому принципу: аппаратное прерывание генерируется контроллером I2C при завершении операции или возникновении ошибки, обработчик верхнего уровня выполняет быструю обработку критически важных операций, а отложенная обработка выполняет дополнительную обработку в контексте процесса.

Такая архитектура обеспечивает минимальное время реакции на критические события и предотвращает потерю данных при высокочастотных операциях I2C. Использование буферизации и очередей сообщений позволяет системе обрабатывать пакетные операции и оптимизировать использование ресурсов шины.

\section{Применение фреймворка}

Разработанный фреймворк был успешно применен в двух крупных промышленных проектах, демонстрируя свою универсальность. Проекты Heating Unit и Midplane представляют различные области применения встраиваемых систем и показывают адаптивность архитектурных решений фреймворка к разнообразным требованиям.

\subsection{Проект Heating Unit}

Проект Heating Unit представляет собой систему управления нагревательными элементами с контролем температуры и влажности. Задача проекта -- обеспечить корректную работу оборудования в условиях низких температур путем предварительного нагрева и поддержания оптимального температурного режима.

Для реализации логики управления Heating Unit был выбран подход на основе конечного автомата (finite state machine). Данный выбор обусловлен следующими факторами:

\begin{itemize}
	\item \textbf{Детерминированное поведение системы} -- конечный автомат обеспечивает предсказуемые реакции на входные воздействия, что критично для систем климат-контроля
	\item \textbf{Четкое разделение режимов работы} -- система управления нагревом естественным образом декомпозируется на отдельные состояния с различной логикой функционирования
	\item \textbf{Наглядность представления алгоритма} -- диаграмма состояний предоставляет интуитивно понятную модель поведения системы
	\item \textbf{Оптимальное использование ресурсов} -- FSM-подход позволяет использовать ограниченные ресурсы микроконтроллера
\end{itemize}

\begin{figure}[H]
	\centering
	\includegraphics[width=\textwidth]{inc/img/struct-diagr2}
	\caption{Машина состояний системы управления нагревом}
	\label{fig:struct_diagr2}
\end{figure}

Представленная на рисунке~\ref{fig:struct_diagr2} схема состояний отражает алгоритм функционирования нагревательного модуля в различных условиях эксплуатации. Система использует четыре основных параметра для определения условий перехода между состояниями:

\begin{itemize}
	\item \textbf{T разр нагрева} -- пороговая температура разрешения активации нагревательного элемента (-5°C по умолчанию)
	\item \textbf{T разр старта} -- минимальная температура окружающей среды для перехода в штатный режим (+5°C по умолчанию)
	\item \textbf{Гистерезис} -- параметр предотвращения частых переключений (2°C по умолчанию)
	\item \textbf{H разр старта} -- пороговое значение влажности для штатной работы (90\% по умолчанию)
\end{itemize}

Диаграмма состояний демонстрирует сложную логику управления, где переходы между состояниями зависят от комбинации температурных и влажностных параметров. Состояние "Инициализация" выполняет первоначальную настройку системы и переходит в соответствующий рабочий режим в зависимости от текущих условий окружающей среды.

Процесс инициализации и загрузки системы представлен в серии блок-схем. Первая диаграмма показывает последовательность загрузки микроконтроллера.

\begin{figure}[H]
	\centering
	\includegraphics[width=0.5\textwidth]{inc/img/block-diagr-hu1}
	\caption{Последовательность загрузки микроконтроллера}
	\label{fig:block_diagr_hu1}
\end{figure}

Схема загрузки в рисунке~\ref{fig:block_diagr_hu1} иллюстрирует критически важный процесс инициализации системы. После подачи питания микроконтроллер выполняет последовательность проверок и настроек, включая тестирование памяти, инициализацию тактовых сигналов и загрузку базовых драйверов. Этот процесс обеспечивает стабильную основу для работы всех последующих компонентов системы.

Следующий этап представляет инициализацию основной программы и общих компонентов системы.

\begin{figure}[H]
	\centering
	\includegraphics[width=0.5\textwidth]{inc/img/block-diagr-hu2}
	\caption{Старт основной программы и общая инициализация}
	\label{fig:block_diagr_hu2}
\end{figure}

Диаграмма на рисунке~\ref{fig:block_diagr_hu2} показывает процесс запуска пользовательского приложения. На этом этапе происходит инициализация фреймворка, настройка подсистем GPIO и I2C, а также подготовка структур данных для работы основного алгоритма. Проверки корректности инициализации обеспечивают надежность системы и предотвращают ошибки во время работы.

Третий этап включает специфичную для платформы инициализацию и запуск основных машин состояний.

\begin{figure}[H]
	\centering
	\includegraphics[width=0.7\textwidth]{inc/img/block-diagr-hu3}
	\caption{Инициализация платформы и основные машины состояний}
	\label{fig:block_diagr_hu3}
\end{figure}

На рисунке~\ref{fig:block_diagr_hu3} представлена схема инициализации платформо-специфичных компонентов и запуска основной логики приложения. Этот этап включает настройку сенсоров температуры и влажности, инициализацию нагревательных элементов и запуск машины состояний климат-контроля. Модульная архитектура фреймворка позволяет легко адаптировать этот процесс к различным аппаратным конфигурациям.

Завершающая диаграмма показывает организацию обработки событий и прерываний в системе.

\begin{figure}[H]
	\centering
	\includegraphics[width=0.5\textwidth]{inc/img/block-diagr-hu4}
	\caption{Обработчики прерываний и таймаутов}
	\label{fig:block_diagr_hu4}
\end{figure}

Схема на рисунке~\ref{fig:block_diagr_hu4} демонстрирует архитектуру обработки асинхронных событий в системе Heating Unit. Система прерываний организована по приоритетам, где критически важные события (например, перегрев) обрабатываются немедленно, а менее критичные события обрабатываются в основном цикле программы. Такой подход обеспечивает быструю реакцию на аварийные ситуации и использование ресурсов процессора.

\subsection{Проект Midplane}

Проект Midplane представляет собой систему управления вентиляторами и мониторинга параметров для серверного оборудования. Задача проекта -- обеспечить корректную работу системы охлаждения путем контроля скорости вращения вентиляторов и мониторинга ключевых параметров системы.

Архитектура проекта Midplane значительно отличается от Heating Unit по сложности и количеству управляемых компонентов. Система включает управление шестью двухроторными вентиляторами через два канала ШИМ с частотой 25 кГц, мониторинг скорости вращения 12 роторов и эмуляцию I2C-устройств для взаимодействия с материнской платой.

\begin{figure}[H]
	\centering
	\includegraphics[width=0.5\textwidth]{inc/img/block-diagr-mid1}
	\caption{Общая архитектура проекта Midplane}
	\label{fig:block_diagr_mid1}
\end{figure}

Общая архитектура Midplane, представленная на рисунке~\ref{fig:block_diagr_mid1}, демонстрирует модульную организацию системы. Центральный блок управления координирует работу всех подсистем, включая контроллеры вентиляторов, систему мониторинга и интерфейсы связи. Такая архитектура обеспечивает масштабируемость системы и упрощает добавление новых функций.

Детализированный процесс инициализации платформы показывает сложность системы Midplane по сравнению с Heating Unit.

\begin{figure}[H]
	\centering
	\includegraphics[width=0.8\textwidth]{inc/img/block-diagr-mid2}
	\caption{Детальная схема инициализации платформы}
	\label{fig:block_diagr_mid2}
\end{figure}

Схема инициализации на рисунке~\ref{fig:block_diagr_mid2} показывает многоэтапный процесс настройки сложной системы. Инициализация включает настройку множественных ШИМ-каналов, конфигурацию счетчиков оборотов для каждого ротора, инициализацию I2C-интерфейсов для связи с материнской платой и настройку системы мониторинга температуры. Каждый этап включает проверки корректности, что обеспечивает надежность системы в производственных условиях.

Система управления интерфейсами реализована через набор специализированных машин состояний.

\begin{figure}[H]
	\centering
	\includegraphics[width=0.9\textwidth]{inc/img/block-diagr-mid3}
	\caption{Машины состояний интерфейсов}
	\label{fig:block_diagr_mid3}
\end{figure}

Диаграмма на рисунке~\ref{fig:block_diagr_mid3} иллюстрирует параллельную работу нескольких машин состояний, каждая из которых отвечает за определенный аспект функционирования системы. Машина состояний управления вентиляторами контролирует скорость вращения в зависимости от температурных показаний, машина состояний мониторинга отслеживает работоспособность всех компонентов, а машина состояний связи обеспечивает взаимодействие с внешними системами через I2C-интерфейс.

Завершающая схема показывает организацию обработки событий в сложной многокомпонентной системе.

\begin{figure}[H]
	\centering
	\includegraphics[width=0.8\textwidth]{inc/img/block-diagr-mid4}
	\caption{Схема обработки событий и прерываний}
	\label{fig:block_diagr_mid4}
\end{figure}

Система обработки событий в проекте Midplane, представленная на рисунке~\ref{fig:block_diagr_mid4}, демонстрирует масштабируемость архитектуры фреймворка. Множественные источники прерываний от различных периферийных устройств обрабатываются через единую диспетчерскую систему, которая распределяет события между соответствующими обработчиками. Приоритизация событий обеспечивает корректную работу критически важных функций системы охлаждения.

Представленные проекты демонстрируют применимость разработанного фреймворка к созданию промышленных решений различной сложности и назначения, подтверждая его универсальность. Модульная архитектура фреймворка обеспечивает гибкость в реализации различных алгоритмов управления, от простых конечных автоматов до сложных многоуровневых систем с множественными интерфейсами. Использование стандартизированных подходов к организации подсистем GPIO и I2C упрощает разработку и сопровождение программного обеспечения, а четкое разделение между аппаратно-зависимыми и аппаратно-независимыми компонентами обеспечивает переносимость решений между различными аппаратными платформами.

Разработанная архитектура фреймворка для микроконтроллера MIK32 обеспечивает решение задач абстракции аппаратных ресурсов и упрощения процесса разработки встроенных приложений.

Многоуровневая структура с четким разделением архитектурно-зависимых, платформо-зависимых и общих компонентов обеспечивает гибкость и переносимость решений. Такой подход позволяет изолировать изменения, связанные с аппаратной платформой, и минимизировать их влияние на прикладные компоненты. Реализованные подсистемы GPIO и I2C демонстрируют принцип архитектурного подхода, основанного на полиморфизме и унификации интерфейсов. Подсистема GPIO обеспечивает единообразную работу с различными типами контроллеров, включая внешние расширители, что подтверждает масштабируемость решения.

Успешная апробация фреймворка в проектах Heating Unit и Midplane различной сложности подтвердила универсальность архитектурного решения и его практическую применимость для промышленных задач. Модульная организация позволяет реализовывать как простые конечные автоматы, так и сложные многокомпонентные системы с параллельными машинами состояний.

Разработанный фреймворк создает техническую основу для ускорения внедрения микроконтроллера MIK32 в промышленные применения и может служить образцом для создания аналогичных решений для других отечественных микроконтроллеров.
